% 本文件由 LaTeX 排版 Agent 根据有效 translation.json 自动生成。
% author=Hippolytus; work_id=0503; title=驳斥犹太人论说

\chapter{驳斥犹太人论说}

% structure: p0001 source=h1 target=omitted reason=page-title-already-rendered-by-parent

1. 现在，请侧耳听我，听我的言语，留心听啊，你这犹太人。你多次自夸，因你定罪了纳匝肋人耶稣致死，给祂喝醋和苦胆；你因此夸耀。所以来吧，让我们一起思考，或许你并非正义地夸耀，以色列啊，那少许的醋和苦胆是否已招致这可畏的警告临到你身上，这是否是你如今身陷无数患难的原因。\par

2. 让那藉着圣神说话并讲真理的——叶瑟的儿子达味——被引到我们面前。他以先知性咏歌唱论真基督，藉圣神赞美我们的天主，（并）清楚地宣告了祂在受难中藉犹太人之手所遭遇的一切；在这咏歌中，那自卑并取了奴仆亚当形像的基督，仿佛以我们的身分呼求天上的天主圣父，在\BibleBook{圣咏}{第六十九篇}中这样说：\BibleQuote{天主，求祢拯救我，因为大水已没入我的灵魂。我深陷在深渊的淤泥中，}这就是说，因乐园中的过犯而陷入阴府的腐朽；\BibleQuote{并且没有立足之处，}就是没有救助。\BibleQuote{我的眼睛因盼望我的天主而失明；（或：因我的盼望）祂几时来拯救我呢？}\par

3. 然后，在接下来的一段中，基督仿佛以自己的身分说：\BibleQuote{那时我偿还了我未曾抢夺的，}这就是说，因亚当的罪，我承受了本非因犯罪而该受的死亡。\BibleQuote{因为，天主啊，祢知道我的愚昧；我的罪过在祢面前不能隐藏，}这是说，\InnerQuote{因为我并未犯罪，}如祂所意指的；为此（加上了）\BibleQuote{愿那些想要看见的人不蒙羞愧}，我的第三天复活，即宗徒们。\BibleQuote{因为为了祢的缘故，}即为了顺服祢，\BibleQuote{我忍受了辱骂，}即十字架，那时\BibleQuote{他们以羞辱蒙蔽我的脸，}即犹太人；那时\BibleQuote{我成了我骨肉兄弟的陌生者，成了我母亲儿女的外人，}（母亲）意指会堂。\BibleQuote{因为为祢殿宇的热忱，父啊，吞灭了我；辱骂祢者的辱骂都落在我身上，}也落在向偶像献祭者身上。因此\BibleQuote{坐在城门口的人议论我，}因为他们把我钉在城门外。\BibleQuote{喝酒的人以歌讥讽我，}即在逾越节筵席上喝酒的人。\BibleQuote{至于我，我在祈祷中向祢，上主说：父啊，求祢宽恕他们，}即外邦人，因为现在是恩待外邦人的时候。\BibleQuote{愿那飓风（试探）不淹没我，深渊（即阴府）不吞灭我：因为祢必不将我的灵魂遗弃在阴府；也不让坑口向我合上，}即坟墓。\BibleQuote{因我的仇敌，求祢解救我，}免得犹太人自夸说：\InnerQuote{我们吞灭了他。}\par

4. 现在基督以人性计划的方式祈求了这一切；然而祂是真天主。但是，如我已说过的，正是那\Quote{奴仆的形像}说了并忍受了这些。因此祂加上：\BibleQuote{我的灵魂期待辱骂和苦难，}这就是说，我甘愿受苦，（而）非被迫。然而\BibleQuote{我等待有人与我同哀，却没有人，}因为我的门徒们都离弃我逃跑了；\BibleQuote{并等待安慰者，我却没有找着。}\par

5. 明智地听啊，犹太人，听听基督说什么：\BibleQuote{他们给了我苦胆作食物；我渴时，他们拿醋给我喝。}这些事祂确实从你那里忍受了。也听听圣神告诉你，祂对那少许醋给了你什么报应。因为先知说，如天主的口吻：\BibleQuote{愿他们的宴席成为网罗和报应。}祂说的是什么报应？显然，是如今抓住你的痛苦。\par

6. 然后听接下来所说的：\BibleQuote{愿他们的眼睛昏黑，以致看不见。}你确实在你的灵魂之眼上，被一种彻底而永久的黑暗所昏黑。因为如今真光已升起，你却如在夜间游荡，在无路之处跌倒，并跌入深渊，因为你离弃了那说\InnerQuote{我是道路}的道路。此外，再听听这句更严重的话：\BibleQuote{并常使他们的脊背弯曲；}这就是说，要使他们给万民作奴隶，不是如同在埃及的四百三十年，也不是如同在巴比伦的七十年，而是说他\Quote{常}使他们屈服于奴役。总之，你如何沉迷于虚妄的希望，指望从抓住你的痛苦中被解救出来？因为那有点奇怪。他向你诅咒这眼睛昏黑，并非不义。而是因为你蒙住了基督的眼睛，（并）如此击打祂，为此，你也常弯曲脊背作奴隶。而当你愤怒地倾倒祂的血时，听听你的报应是什么：\BibleQuote{求祢倾泄祢的怒愤在他们身上，愿祢的怒怒抓住他们；}以及\BibleQuote{愿他们的住所荒凉，}即他们著名的圣殿。\par

7. 但是，先知啊，请告诉我们，圣殿为何变为荒凉？是因那古时铸造牛犊的事吗？是因百姓崇拜偶像吗？是因先知的血吗？是因以色列的奸淫和淫乱吗？决不是，他说；因为在这一切过犯中，他们总能找到敞开的宽恕和仁慈；而是因为他们杀害了他们恩主的圣子（祂与圣父同永恒）。因此祂说：\Quote{父啊，使他们的殿变为荒凉；因为他们迫害了祢按自己旨意为了世界的救恩所击打的那一位；}即：他们以暴力和不义的死亡迫害了我，\Quote{并加增了我伤口的痛苦。}从前，作为热爱人类者，我为外邦人的迷失而痛苦；但在这痛苦之上，他们又加上了另一层，因他们自己也迷失了。因此\Quote{将罪孽加于他们的罪孽，将患难加于患难，不要让他们进入祢的正义，}即进入祢的国度；但是\Quote{让他们从生命册上被涂抹，不得与义人一同被记录，}即与他们的圣祖和族长一同被记录。\par

8. 对此你有何话说，犹太人？说这些话的既不是玛窦，也不是保禄，而是你的受傅者达味，他因基督的缘故宣告并宣布了这些可怕的判决。如同伟大的约伯一样，他对那反对正义和真实者的你说道：\Quote{你像出卖奴隶一样出卖了基督，你像强盗一样在园中来到祂面前。}\par

9. 我现在提出撒罗满关于基督的预言，它清楚明白地宣告了关于犹太人的事；以及那些不仅现在临到他们，而且将来在未来的世代也要临到他们的事，都是因为他们对生命之君所表现的顽梗和狂妄；因为先知说：\Quote{不敬虔的人说，彼此议论，却不正确，}即关于基督，\Quote{我们埋伏等候义人，因为他与我们不合，他完全反对我们的言行，指责我们违反法律，并自称认识天主；他称自己为天主的儿子。}然后他说：\Quote{他连看我们也是可厌的；因为他的生命不像别人，他的道路是另一种样式。我们被他视为伪君子，他回避我们的道路如同回避污秽，并宣告义人的结局是有福的。}再者，犹太人！请听这话：没有哪位义人或先知曾称自己为天主子。因此，撒罗满再次以犹太人的身份论及这位义者，即基督，这样说道：\Quote{他被造是为了责备我们的意念，并夸耀天主是他的父。让我们看看他的话是否真实，让我们试验他结局如何；因为如果那义人是天主子，天主必帮助他，救他脱离仇敌的手。让我们用羞辱的死来定他的罪，因为照他自己的话，他必被尊重。}\par

10. 再者，达味在圣咏中论到未来的世代说：\Quote{那时祂}（即基督）\Quote{要在祂的震怒中对他们说话，在祂的烈怒中使他们苦恼。}再者，撒罗满论到基督和犹太人时说：\Quote{当义人站在那些曾压迫他、不理会他话的人面前，大有胆量时，他们看见就必在极大的恐惧中惊慌，并因他救恩的奇异而惊奇；他们必后悔，心中痛苦叹息，彼此说：这就是我们曾一度讥笑、当作笑柄的人；我们愚昧人以为他的生命是疯狂，他的结局是没有荣耀。他怎么被列在天主的儿女中，他的分在圣徒中呢？因此我们偏离了真理之道，正义之光没有照耀我们，正义的太阳也没有升起在我们身上。我们在邪恶和灭亡的道路上疲乏；我们走过没有道路的旷野；至于上主的道路，我们却不曾知道。我们的骄傲有什么益处呢？这一切都像影子一样消逝了。}\par

结论缺失。\par

\SourceRecord{Translated by J.H. MacMahon.；Ante-Nicene Fathers；Vol. 5.；Edited by Alexander Roberts, James Donaldson, and A. Cleveland Coxe.；Buffalo, NY: Christian Literature Publishing Co.,；1886.；<http://www.newadvent.org/fathers/0503.htm>.}
